\documentclass[11pt,a4paper]{article}
\usepackage[margin=1in]{geometry}
\usepackage{enumitem}
\usepackage{titlesec}
\usepackage{hyperref}
\usepackage{parskip}

\hypersetup{
    colorlinks=true,
    linkcolor=blue,
    urlcolor=blue,
    citecolor=blue
}

\titleformat{\section}{\large\bfseries}{\thesection}{1em}{}[\titlerule]
\titleformat{\subsection}{\normalsize\bfseries}{\thesubsection}{1em}{}

\pagestyle{plain}

\begin{document}

\begin{center}
{\LARGE\textbf{Curriculum Vitae}}\\[0.3cm]
{\Large Daniel Rogozin}
\end{center}

\section{Education}

\begin{itemize}[leftmargin=*]
    \item PhD programme at the Institute for Information Transmission Problems (2020--2022) under supervision by Ilya Shapirovsky (unfinished because of the sudden immigration in March 2022 caused by the political reasons).
    \item Master's degree at the Lomonosov Moscow State University, The Faculty of Philosophy (2016--2018).
    \item Bachelor's degree at the State Academic University for the Humanities, The Faculty of Philosophy (2012--2016).
\end{itemize}

\section{Theses}

\begin{itemize}[leftmargin=*]
    \item 2018, \textit{Modal type theory based on intuitionistic epistemic logic}, Master's thesis, Lomonosov Moscow State University.
    \item 2016, \textit{Type theory in the foundations of mathematics}, Bachelor's thesis, State Academic University for the Humanities.
\end{itemize}

\section{Employment}

\begin{itemize}[leftmargin=*]
    \item (March 26 ---) Corca Research, Mathematician.
    \item (Sep' 23---Feb' 26) Noeon Research, Research Scientist. 
    
    My research at Noeon was about mathematical and philosophical conceptualisation of engineering ideas and transforming practice-related requirements
    into a coherent mathematical theory. I contributed to a number of internal projects related to type theory and category theory with applications to interpretable
    artificial intelligence and knowledge representation:
\begin{itemize}
    \item I suggested a proof-theoretic conceptualisation of the category theory-based reasoning framework that resulted in a mathematical preprint ``Term Assignment and Categorical Models for Intuitionistic Linear Logic with Subexponentials''
    \item Developing a Haskell type checker for the DreamCoder-based program synthesis utility.
    \item Permanently reviewing amounts of the current literature in programming language theory,
    \item Educate colleagues in mathematical logic, philosophy of language, denotational semantics, universal algebra, type theory, programming language theory, etc.
\end{itemize}
    \item (Jan' 23---June 23) University College London, Department of Computer Science, teaching assistant.
    \item (May' 22---April' 23) Yatima Inc/Lurk Lab, Lean 4 engineer. The projects include:
    \begin{itemize}
    \item The Lean 4 formalisation of finite fields, elliptic curve cryptography and the zero-knowledge proof generation framework Nova in Lean.
    \end{itemize} 
    \item (Feb' 18---Jun' 22) Serokell OÜ, Tallinn, Estonia, Senior Haskell software engineer. The projects include:
    \begin{itemize}[leftmargin=*]
    \item (2018) Cardano, a cryptocurrency backend in Haskell.
    \item (2018) Ariadne, an R\&D project, a cryptocurrency wallet.
    \item (2018) Snowdrop, an R\&D project, a framework for blockchain systems with a proof-of-stake consensus.
    \item (2019) Machine Learning in Haskell, an R\&D project on type-safe matrix operations via singletons. The work on that project also included the preparation of blogposts and the conference talk.
    \item (2019) Formality, an R\&D project, a dependently typed language with affine types.
    \item (2020--2022) Glasgow Haskell Compiler contribution: the implementation of the Char Kind and partial implementation of @-binders in function declarations.
    \end{itemize}
    \item (Jan' 21---Mar' 22) Institute for Information Transmission Problems of Russian Academy of Sciences, Laboratory of Algebra and Number Theory, Trainee researcher.
\end{itemize}

\section{Research Areas}

\begin{itemize}[leftmargin=*]
    \item \textbf{Modal logic}: semantic analysis of constructive modal logics and decidability aspects of classical modal logics.
    \item \textbf{Categorical logic}: category-theoretic models of intuitionistic modal and linear logics and Kripke-Joyal semantics.
    \item \textbf{Algebraic logic}: representation theory for Boolean algebras with operators.
    \item \textbf{Programming language theory}: functional languages, formal verification with dependent types.
    \item \textbf{Knowledge Representation}: mathematical and philosophical foundations of knowledge representation as a theoretical discipline.
\end{itemize}

\section{Publications}

\begin{itemize}[leftmargin=*]
    \item Daniel Rogozin. ``Semantic Analysis of Subexponential Modalities in Distributive Non-commutative Linear Logic'', \emph{Proceedings Modalities in substructural logics: Applications at the interfaces of logic, language and computation (AMSLO 2023)}, \emph{Electronic Proceedings in Theoretical Computer Science}, 381, pp. 60–70.
    \item Daniel Rogozin and Ilya Shapirovsky. ``On decidable extensions of Propositional Dynamic Logic'', \emph{Journal of Applied Logics – IfCoLog Journal}, 10(4): 629–642, 2023.
    \item Amin, Nada, John Burnham, François Garillot, Rosario Gennaro, Daniel Rogozin, Chhi'med Kunzang and Cameron Wong. ``Lurk: Lambda, the ultimate recursive knowledge.'' \emph{Proceedings of the ACM on Programming Languages}, 7(ICFP), 259-274.
    \item Daniel Rogozin. ``Some results on relation algebra reducts: Residuated and semilattice-ordered semigroups.'' \textit{Journal of Logic and Computation} 32, (2022): 266--280.
    \item Daniel Rogozin. ``Reducts of relation algebras: The aspects of axiomatisability and finite representability.'' In \textit{International Symposium on Logical Foundations of Computer Science}, pp. 266--280. Springer, Cham, 2022.
    \item Daniel Rogozin. ``Categorical and algebraic aspects of the intuitionistic modal logic IEL$^-$ and its predicate extensions.'' \textit{Journal of Logic and Computation} 31, no. 1 (2021): 347--374.
    \item Daniel Rogozin. ``Modal Type Theory Based on the Intuitionistic Modal Logic IEL$^-$.'' In \textit{International Symposium on Logical Foundations of Computer Science}, pp. 236--248. Springer, Cham, 2020.
\end{itemize}

\section{Preprints}

\begin{itemize}[leftmargin=*]
    \item Rogozin, Daniel. ``Term Assignment and Categorical Models for Intuitionistic Linear Logic with Subexponentials.'' arXiv preprint arXiv:2507.12360 (2025).
    \item Daniel Rogozin. ``The distributive full Lambek calculus with modal operators.'' arXiv preprint arXiv:2003.09975 (2020).
    \item Daniel Rogozin. ``Quantale semantics of Lambek calculus with subexponential modalities.'' arXiv preprint arXiv:1908.01055 (2019).
\end{itemize}

\section{Conference Talks}

\begin{itemize}[leftmargin=*]
    \item Winter 2026, workshop ``Formalisation and deformalisation of mathematical reasoning'', Paris, France, \emph{Operational Knowledge Representation with Monoidal Categories and Linear Type Theory}.
    \item September 2025, conference ``British Logic Colloquium 2025'', Manchester, UK, \textit{Term Assignment and Categorical Models for Intuitionistic Linear Logic with Subexponentials}.
    \item July 2023, 34th European Summer School in Logic, Language and Information, workshop ``Modalities in substructural logics: applications at the interfaces of logic, language and computation'', Ljubljana, Slovenia, \textit{Semantic Analysis of Subexponential Modalities in Distributive Non-commutative Linear Logic}.
    \item August 2022, conference ``Advances in Modal Logic 2022'', Rennes, France, \textit{On decidable extensions of Propositional Dynamic Logic} (joint work with Ilya Shapirovsky).
    \item January 2022, conference ``Logical Foundations of Computer Science 2022'', Deerfield Beach, Florida, USA, \textit{Reducts of Relation Algebras: The Aspects of Axiomatisability and Finite Representability}.
    \item June 2021, conference ``BLAST: Boolean Algebras, Lattices, Algebraic and Quantum logic, Universal Algebra, Set Theory Set-theoretic and Point-free Topology'', New Mexico State University, New Mexico, USA, \textit{Reducts of representable relation algebras: finite representability and axiomatisability}.
    \item June 2021, workshop ``Logical Perspectives 2021'', Steklov Institute of Mathematics of Russian Academy of Sciences, Moscow, Russia, \textit{The distributive full Lambek calculus with modal operators}.
    \item August 2020, conference ``Advances in Modal Logic'', Helsinki, Finland, \textit{Canonical Extensions for the Distributive Full Lambek Calculus with Modal Operators}.
    \item January 2020, conference ``Logical Foundations of Computer Science 2022'', Deerfield Beach, Florida, USA, \textit{Modal Type Theory Based on the Intuitionistic Modal Logic ${\bf IEL}^{-}$}.
    \item December 2019, workshop ``Logic matters'', Higher School of Economics, The Faculty of Computer Science, \textit{Semantics of the Lambek calculus and its modal extensions: approaches and open issues}.
    \item June 2019, conference ``Topology, Algebra, and Categories in Logic'', Nice, France, \textit{Quantale semantics for Lambek calculus with subexponentials}.
\end{itemize}

\section{Seminar Talks}

\begin{itemize}[leftmargin=*]
    \item Winter 2023, PPLV group seminar, UCL, London, \textit{Intuitionistic epistemic logic categorically and algebraically}.
    \item Winter 2023, PPLV group seminar, UCL, London, \textit{Decidable extensions of propositional dynamic logic}.
    \item Spring 2022, New Mexico State University, New Mexico, \textit{Categories and lambda calculus for programming language theory}.
    \item Winter 2022, The Seminar on Modal and Algebraic Logic, Steklov Mathematical Institute of Russian Academy of Sciences, \textit{Relation algebras: overview and open questions}.
    \item Autumn 2021, seminar ``Computational Logic'', City University of New York, The Graduate Centre, \textit{Decidability results for relation algebras and their reducts}.
    \item Autumn 2021, seminar ``Constructive Knowledge'', Institute of Philosophy of Russian Academy of Sciences, Moscow, \textit{A version of Kripke-Joyal semantics for intuitionistic modal predicate logics}.
    \item Autumn 2021, seminar ``From the Logical Point of View'', Higher School of Economics, The Faculty of Philosophy, \textit{A survey of relation algebras}.
    \item Spring 2021, The Seminar on Modal and Algebraic Logic, Lomonosov Moscow State University, The Faculty of Mechanics and Mathematics, \textit{Representation of cylindric algebras with back-and-forth games}.
    \item Autumn 2020, seminar ``Computational Logic'', City University of New York, The Graduate Centre, \textit{Categorical and Algebraic Aspects of the Intuitionistic Modal Logic IEL$^-$ and its predicate extensions}.
    \item Spring 2020, The Logical Seminar of the Laboratory of Mathematical logic, St. Petersburg Department of Mathematical Institute of Russian Academy of Sciences, St. Petersburg, \textit{Canonical extensions for modal logics}.
    \item Spring 2020, The Seminar on Modal and Algebraic Logic, Lomonosov Moscow State University, The Faculty of Mechanics and Mathematics, \textit{Topological duality and canonical extensions for logics of distributive residuated lattices with modal operators}.
    \item Winter 2020, seminar ``Contemporary problems of mathematical logic'', Higher School of Economics, The Faculty of Mathematics, \textit{Introduction to Canonical extensions}.
    \item Autumn 2019, seminar ``Constructive Knowledge'', Institute of Philosophy of Russian Academy of Sciences, Moscow, \textit{Concept analysis via substructural logics}.
    \item Autumn 2019, The Seminar on Modal and Algebraic Logic, Lomonosov Moscow State University, The Faculty of Mechanics and Mathematics, \textit{Topological duality for classical modal logics}.
    \item Autumn 2019, seminar ``Contemporary problems of mathematical logic'' seminar, Higher School of Economics, The Faculty of Mathematics, \textit{Duality for Heyting algebras}.
    \item Autumn 2019, The Proof Theory seminar, Steklov Institute of Mathematics of Russian Academy of Sciences, Moscow, \textit{Discussion on the paper ``Locales, nuclei, and Dragalin frames'' by G. Bezhanishvili and W. Holliday}.
    \item Spring 2019, seminar ``Constructive Knowledge'', Institute of Philosophy of Russian Academy of Sciences, Moscow, \textit{Formal concept lattices}.
    \item Spring 2019, seminar ``Constructive Knowledge'', Institute of Philosophy of Russian Academy of Sciences, Moscow, \textit{The univalence axiom in Coq}.
    \item Spring 2019, The Proof Theory seminar, Steklov Institute of Mathematics of Russian Academy of Sciences, Moscow, \textit{Semantics of the Lambek calculus with subexponentials}.
    \item Autumn 2018, seminar ``Constructive Knowledge'', Institute of Philosophy of Russian Academy of Sciences, Moscow, \textit{Formal verification for blockchain}.
    \item Spring 2018, The Constructive Knowledge seminar, Institute of Philosophy of Russian Academy of Sciences, Moscow, \textit{Categorical semantics of the type theory based on the Intuitionistic Epistemic Logic}.
    \item Spring 2018, The seminar on non-classical logics and computability, Lomonosov Moscow State University, Moscow, \textit{Type theory for the intuitionistic modal logic IEL$^-$}.
    \item Autumn 2017, The Constructive Knowledge seminar, Institute of Philosophy of Russian Academy of Sciences, Moscow, \textit{Modal Type Theory based on the Intuitionistic Epistemic Logic}.
    \item Spring 2017, The Constructive Knowledge seminar, Institute of Philosophy of Russian Academy of Sciences, Moscow, \textit{Curry-Howard correspondence and Kolmogorov complexity}.
    \item Autumn 2016, The Constructive Knowledge seminar, Institute of Philosophy of Russian Academy of Sciences, Moscow, \textit{Type theory, intuitionistic logic, and Kolmogorov complexity}.
    \item Spring 2016, The Constructive Knowledge seminar, Institute of Philosophy of Russian Academy of Sciences, Moscow, \textit{Type theory and proof formalisation}.
\end{itemize}

\section{Research Projects}

I participated in the following projects as a collaborator.
\begin{itemize}[leftmargin=*]
    \item The research grant MK-1184.2021.1.1 (2021--2022), \textit{Algorithmic and semantic questions of modal extensions of linear logics}.
    \item The research grant MK-430.2019.1 (2019--2021), \textit{Subexponential modalities in noncommutative linear logic}.
    \item The research RFBR grant 19-011-00799 (2019--2021), \textit{Knowledge Justification in Formal Epistemology}.
    \item The research RFBR grant 16-03-00364 (2016--2018), \textit{Logical and Epistemological Aspects of Constructive Knowledge}.
\end{itemize}

\section{Grants}

\begin{itemize}[leftmargin=*]
    \item London Mathematical Society Solidarity Supplementary Grant, (2022--2023).
    \item Isaac Newton Institute Solidarity Scheme, (2022--2023).
\end{itemize}

\section{Teaching}

\begin{itemize}[leftmargin=*]
    \item Winter and Spring 2023, University College London, The Department of Computer Science, London, UK, \textit{Introduction to Logic and Algorithms}, teaching assistant.
    \item Autumn 2021, Higher School of Economics, The Faculty of Computer Science, Moscow, \textit{Functional programming} (with Stepan Kuznetsov), lecturer and teaching assistant.
    \item Autumn 2021, ITMO University, The Department of CS, St. Petersburg, \textit{Functional Programming}, lecturer and teaching assistant.
    \item Autumn 2021, Computer Science Club, St. Petersburg, \textit{Semantic aspects of intuitionistic logic}, lecturer.
    \item Summer 2021, Lalalambda summer school, Moscow, \textit{Introduction to modal and temporal logics}, lecturer.
    \item Spring 2021, Higher School of Economics, The Faculty of Computer Science, Moscow, \textit{Functional programming} (with Stepan Kuznetsov), lecturer and teaching assistant.
    \item Spring 2020, Higher School of Economics, The Faculty of Computer Science, Moscow, \textit{Functional programming} (with Stepan Kuznetsov), lecturer and teaching assistant.
    \item Spring 2020, ITMO University, The Department of CS, St. Petersburg, \textit{Functional Programming}, lecturer and teaching assistant.
    \item Autumn 2019, Lomonosov Moscow State University, \textit{Introduction to Logic}, teaching assistant.
    \item Autumn 2019, Computer Science Club, St. Petersburg, \textit{Introduction to modal logic}, lecturer.
    \item Summer 2019, Computer Summer School, Narva, lecture \textit{Introduction to computability and logic}.
    \item Spring 2019, Mathematical Class at Moscow Pedagogical State University, lecture \textit{Proofs and computations}.
    \item Spring 2019, Higher School of Economics, The Faculty of Computer Science, Moscow, \textit{Functional programming} (with Stepan Kuznetsov), lecturer and teaching assistant.
    \item Spring 2019, ITMO University, The Department of CS, St. Petersburg, \textit{Functional Programming}, lecturer and teaching assistant.
    \item Autumn 2018, Lomonosov Moscow State University, \textit{Introduction to Logic}, teaching assistant.
\end{itemize}

\section{Paper Reviewing}

I was a reviewer for the conferences ``Advances in Modal Logic 2020'', ``The Eighth International Conference On Logic, Rationality, and Interaction'', ``International Colloquium on Automata, Languages and Programming 2023'', ``31st EACSL Annual Conference on Computer Science Logic'',
``Relational and Algebraic Methods in Computer Science''.

\section{Blogposts and YouTube Videos}

\begin{itemize}[leftmargin=*]
    \item ``A Walk Through L: Operationalising Knowing How with Category Theory and Linear Logic'', \url{https://subexponentials.substack.com/p/a-walk-through-l}
    \item ``Algebraic data types in 10 minutes. Haskell Tutorial'', \url{https://youtu.be/UqwLn2OyQ_E}
    \item ``Operators in 10 minutes. Haskell Tutorial'', \url{https://youtu.be/O1K8EMiPEVE}
    \item ``Philosophy and metamathematics'', \url{https://youtu.be/GjBgWk83GVw}
    \item ``How Dependent Haskell Can Improve Industry Projects'' (written with Vladislav Zavialov), \url{https://serokell.io/blog/how-dependent-haskell-can-improve-industry-projects}
    \item ``Formal verification: History and Methods'', \url{https://serokell.io/blog/formal-verification-history}
    \item ``Incomplete and Utter Introduction to Modal Logic, Pt. 1'', \url{https://serokell.io/blog/incomplete-and-utter-introduction-to-modal-logic}
    \item ``Incomplete and Utter Introduction to Modal Logic, Pt. 2'', \url{https://serokell.io/blog/rapid-introduction-to-modal-logic-2}
    \item ``Dimensions and Haskell: Introduction'', \url{https://serokell.io/blog/dimensions-and-haskell-introduction}
    \item ``Dimensions and Haskell: Singletons in Action'', \url{https://serokell.io/blog/dimensions-haskell-singletons}
    \item ``Constructive and Non-Constructive Proofs in Agda (Part 1): Logical Background'', \url{https://serokell.io/blog/logical-background}
    \item ``Constructive and Non-Constructive Proofs in Agda (Part 2): Agda in a Nutshell'', \url{https://serokell.io/blog/agda-in-nutshell}
    \item ``Constructive and Non-Constructive Proofs in Agda (Part 3): Playing with Negation'', \url{https://serokell.io/blog/playing-with-negation}
\end{itemize}

\section{Awards}

\begin{itemize}[leftmargin=*]
    \item 2022, Best student paper award, for the paper \textit{Reducts of relation algebras: The aspects of axiomatisability and finite representability} at Logical Foundations of CS 2022, Miami.
    \item 2021, Category ``Best Course for New Knowledge and Skills'' for the Functional programming course at Higher School of Economics, Moscow.
    \item 2020, Category ``Best Course for New Knowledge and Skills'' for the Functional programming course at Higher School of Economics, Moscow.
\end{itemize}

\section{Languages}

\begin{itemize}
    \item Russian (native speaker),
    \item English (fluent, close to native),
    \item French (actively learning, currently ~A2),
    \item Yiddish (learning for cultural curiousity),
\end{itemize}

\section{Contact Info and Links}

\begin{itemize}[leftmargin=*]
    \item Phone number: +447438533004
    \item Email: \href{mailto:relationalgebra@gmail.com}{relationalgebra@gmail.com}
    \item Location: London, United Kingdom
    \item Github: \url{https://github.com/DanielRrr}
    \item LinkedIn: \url{https://www.linkedin.com/in/daniel-rogozin-a94659265/}
    \item Google Scholar: \url{https://scholar.google.com/citations?user=t3Vf1oAAAAAJ&hl=en}
\end{itemize}

\end{document}